\chapter{为什么地址不能直接等同于物理位置}

\begin{statebox}
\textbf{项目里程碑 v0.3。} Stage 1 已经位于低端 RAM，但处理器仍不能执行 64 位内核。
它先用 PML4、PDPT 和 PD 恒等映射低 64 MiB，再按顺序提交 PAE、CR3、LME 与分页。
这份启动页表解决“怎样进入长模式”，却不能解决多个程序使用相同地址、代码与数据权限
不同以及物理页需要动态归属的问题。
\end{statebox}

\inputchapterbody{chapters/ch11-virtual-memory}

\section{从缺页现场到可测试地址空间}
\inputdetail{chapters/ch11-stage-three-cpp20-memory-model}
\inputdetail{chapters/ch00-stage-one-page-fault-source-walk}

\begin{problemframe}
\textbf{尚未解决的问题。} 地址已经能够由页表解释，映射也有权限、缺页和回收规则，但
磁盘中的字节仍不是一份可进入的内核映像。下一章回到项目 v0.4，验证 Kernel 容器和
ELF64 程序头，在不留下半提交段的条件下建立内核并交接 BootInfo。
\end{problemframe}
